\begin{abstract}
This is Version~8 of the scoped O.1 cone audit; it revises
XVI-O1-CONE-v7. The cone \(C_\alpha\times\RR^2\), its deficit and
holonomy, and the curvature-regular (CR) and componentwise
Friedrichs (F) realizations remain prescribed inputs.
The v7 analytical content is retained. What is new is executed
and is stated with its scope.

\emph{(1) The v7 program was run.} The embedded \loc{cone_v7.py}
was extracted unchanged and executed. Its self-test passed, all
27 replay runs returned 0, and the coverage-aware comparison with
the historical logs passed for all \(18\) graph--branch pairs.
Of the fifteen static corrections listed in v7, eight are real:
seven code defects and one methodological limitation. None of them changes a number printed in v6.

\emph{(2) The count of 26 is now exact, and so is a complete span
certificate.} We evaluated the apex functionals on every
relevant unit atom over a complete \(6\times6\) momentum grid,
in exact \(\mathbb Q(\ii)\) arithmetic. At
\(\alpha=3/10\) (CR and F) and at \(\alpha=7071/10000\) (CR), the
inventory is \(34=18+8+6+2\). The Slavnov--Taylor row rank is
\(26\), with a nonzero minor and a kernel verified on all rows.
The closed space is \(Z\) with \(\dim Z=8\). The exact directions
\(B\) satisfy \(\dim B=6\) and \(B\subset Z\), and
\(\dim H^0=2\). At \(\alpha=7071/10000\) the F branch is different:
both the historical modular record and the v8 exact run give
\(38=21+8+7+2\), with ST rank \(29\), \(\dim Z=9\) and
\(\dim H^0=3\).
The number 34 is therefore a branch- and angle-qualified count,
not a universal one.

\emph{(3) Projections.} Six of the 26 Ward combinations are
supported on the twelve constant-profile coordinates.
Two of them lie in the \(hh\) block: \(b_1+2b_3=0\) and \(b_2=0\).
Both hold exactly on the \(hh\) catalogue sum. They hold through
a cancellation between D01 and D03; neither graph satisfies them
alone. The other four involve the CK/CB sectors, and we identify
them with the v7 rows only conditionally. The remaining 20 are
open.
Of the 18 engine \(hh\) coordinates, the six constant-profile
ones receive an explicit exact dictionary; the other twelve stay
open.
Of the six exact values, four combinations are fixed, including
\(e_1=(1-\alpha)/4\) and \(e_3=0\). The CR/F difference sits only
in \(e_4+e_5\), where it equals \(\nu/8\) exactly at both
computed angles. The remaining two
combinations are open.

\emph{(4) New mixed-channel finding.} The engine slice exposes
hidden mixed pairs. These are pairs whose circle pole vanishes
while their Feynman density does not. In D01, the lattice keys
\((4,4),(4,7),(5,5),(7,4),(7,7)\) carry a mixed
\((0,2)\) density proportional to \(P^2(6x^2-10x+3)\).
The printed v6 \(hh\) branch difference is therefore unchanged in
value. Its status becomes conditional on the corresponding
mixed-channel trace. A per-mode Duhamel identity and a reduction
formula are proved, but the D04 part of the second \(\phi\phi\)
value stays open.

\emph{(5) Evanescent coefficients.} Under an explicitly declared
slice-DRED prescription with formal \(\tilde g\) components, the
number of evanescent pole coefficients is 86 on each physical
branch and 113 on the UNIT lattice. They are prescription
conditional.

No missing value is set to zero.
The exact rank statement has status
\texttt{CLOSED\_PENDING\_}\allowbreak\texttt{INDEPENDENT\_REVIEW}
only on its precisely stated finite scope. Overall:
\[
 \text{O.1 on the cone}=\NOTCLOSED,\qquad
 \texttt{bulk\_symmetry\_restored}=\texttt{false}.
\]
\end{abstract}
